\documentclass[11pt]{article}

\usepackage[margin=2.5cm]{geometry}
\usepackage{amsmath,amssymb}
\usepackage{booktabs}
\usepackage{array}
\usepackage{graphicx}
\usepackage{longtable}
\usepackage{pdflscape}
\usepackage{microtype}
\usepackage{xcolor}
\usepackage{xspace}
\usepackage{hyperref}

\hypersetup{
  colorlinks=true,
  linkcolor=blue!55!black,
  citecolor=blue!55!black,
  urlcolor=blue!55!black
}

\newcommand{\IPPL}{\textnormal{\textsc{ippl}}\xspace}
\newcommand{\todo}[1]{\textcolor{red!70!black}{[TODO: #1]}}
\newcommand{\validationplot}[3]{%
  \begin{landscape}
  \begin{figure}[p]
    \centering
    \includegraphics[width=\linewidth,height=.78\textheight,keepaspectratio]{figures/validation/#1}
    \caption{#2}\label{#3}
  \end{figure}
  \end{landscape}%
}

\title{A Performance-Portable Particle--Mesh Framework for Cosmological
Structure Formation with \IPPL}
\author{\todo{author list and affiliations}}
\date{Draft of \today}

\begin{document}
\maketitle

\begin{abstract}
We present the first cosmological structure-formation application built on the
Independent Parallel Particle Layer (\IPPL), a C++ framework for distributed,
performance-portable particle--mesh algorithms.  The present model evolves one
collisionless cold-matter species in a periodic flat $\Lambda$CDM universe.  It
combines Gaussian first-order Lagrangian perturbation theory initial conditions,
cloud-in-cell particle--mesh coupling, a distributed spectral Poisson solve, and
kick--drift--kick integration in the scale factor.  The implementation inherits
the MPI, Kokkos, and heFFTe abstractions previously exercised by \IPPL's plasma
particle-in-cell mini-applications, while adding cosmological background
evolution, transfer functions, reproducible Fourier-space initialization, and
cosmology-specific validation.

The validation establishes linear evolution on one to four MPI ranks with an
OpenMP backend and matched initial-condition statistics against an independent
generator.  It also tests frozen forces, analytical pre-crossing trajectories,
and nonlinear evolution against native plain-PM FastPM.  Matched-discretization
agreement is close, but crossed particle/mesh studies retain resolution failures;
the completed Gaussian matrix does not yet establish a qualified resolution
band.  These bounded results position the code as a portable particle--mesh
research and verification platform, not a replacement for mature production
cosmology codes.  Higher-resolution studies, second-order initial conditions,
short-range gravity, scalable output, and accelerator and multi-node qualification
remain necessary for broader scientific and performance claims.
\end{abstract}

\paragraph{Scope and contribution.}
The contribution is twofold.  First, it translates a validated particle--mesh
software stack from plasma physics to cosmological gravity without duplicating
the portability layer.  Second, it defines a staged validation path in which
background cosmology, initial conditions, discretized force, time integration,
parallel invariance, and performance are tested independently before nonlinear
science claims are made.  The initial linear results were summarized at
development snapshot \texttt{0f2953bde}; Appendix~\ref{app:validation} records
the completed validation campaigns through 4 October 2026, their separate
provenance, and their retained failures.  Results marked \todo{measurement}
are planned rather than completed.

\subsection*{Position in computational cosmology}

Cosmological simulation is a mature field with several distinct algorithmic
families.  GADGET-4 combines tree or fast-multipole short-range gravity with a
particle--mesh long-range solver and includes hydrodynamics, on-the-fly structure
finding, and second-order initial conditions~\cite{springel2021gadget4}.  SWIFT
uses fine-grained tasking and couples fast-multipole gravity to an optional
Fourier particle--mesh solver, alongside hydrodynamics and galaxy-formation
models~\cite{schaller2024swift}.  HACC was designed for extreme-scale
cosmological simulations and uses an architecture-dependent short-range solver
around a common long-range spectral method~\cite{habib2016hacc}.  Abacus attains
high force accuracy and throughput with a specialized near/far decomposition;
the AbacusSummit campaign evolved roughly 60 trillion particles over many
cosmologies~\cite{maksimova2021abacussummit}.  Mesh and adaptive-mesh codes such
as RAMSES and Nyx occupy a complementary Eulerian niche, particularly when gas
dynamics and refinement are central~\cite{teyssier2002ramses,almgren2013nyx}.

Pure and approximate particle--mesh methods are an equally important comparison.
FastPM modifies the kick and drift operators to preserve linear displacement
growth with few time steps and targets fast mock-catalog generation, while COLA
evolves residual motion around perturbative trajectories for a similar
purpose~\cite{feng2016fastpm,tassev2013cola}.  CONCEPT couples a PM long-range
solver to a short-range particle--particle method and adaptive stepping, with
tight integration of cosmological initial conditions and analysis
products~\cite{dakin2022concept}.  Thus neither particle--mesh cosmology nor
high-throughput ensemble generation is by itself a novelty claim for this work.
The research question is instead whether a reusable particle-and-field framework
can retain cosmological accuracy and useful efficiency while holding the
high-level implementation fixed across vendor architectures.

Against these production systems, the near-term role of an \IPPL-based code is
not feature parity.  Its distinctive value is a compact, open particle--mesh
test bed in which the same application-level C++ code targets multiple CPU and
GPU programming backends through Kokkos, distributed FFTs through heFFTe, and
domain decomposition through \IPPL.  This makes it suitable for controlled
studies of portable deposition and interpolation kernels, FFT-based gravity,
reproducible distributed initialization, load balancing, mixed precision, and
new Poisson or time-integration algorithms.  The approach mirrors the role of
the ALPINE plasma mini-apps as an interface between physics, algorithms, and
computer architecture~\cite{muralikrishnan2024scaling}.

Table~\ref{tab:positioning} summarizes the intended positioning.  The immediate
scientific target is periodic collisionless large-scale structure in the
linear and mildly nonlinear regimes.  High-dynamic-range halo and galaxy
formation requires a short-range force method and, for baryonic science,
hydrodynamics and subgrid physics.  Those additions should be motivated by a
specific science case rather than treated as prerequisites for a useful
particle--mesh research code.

\begin{table}[htbp]
  \centering
  \caption{A functional comparison that clarifies the intended role of the
  \IPPL cosmology application.  ``Planned'' denotes work not qualified by the
  results in this paper.}
  \label{tab:positioning}
  \begin{tabular}{@{}p{0.23\linewidth}p{0.31\linewidth}p{0.36\linewidth}@{}}
    \toprule
    Capability & Mature production landscape & \IPPL cosmology status and role \\
    \midrule
    Periodic long-range gravity & Spectral PM is common in HACC, TreePM/FMM--PM
      codes, and other large-volume solvers & Implemented with CIC and a
      distributed spectral Poisson solve; primary current capability \\
    Short-range force & Trees, FMM, P$^3$M, AMR, or specialized near-field
      algorithms & Not implemented; required for halo-scale accuracy \\
    Initial conditions & 2LPT and external production generators are standard &
      Deterministic distributed 1LPT, BBKS and tabulated transfer functions;
      2LPT is planned \\
    Heterogeneous hardware & Often optimized through code-specific CPU/GPU
      paths & Single-source Kokkos abstraction inherited from \IPPL; accelerator
      cosmology runs remain to be qualified \\
    Baryonic physics & Available in codes such as SWIFT, GADGET-4, Nyx, and
      RAMSES & Outside the present paper; possible through future IPPL fields or
      coupling to a hydrodynamics package \\
    Research use & Precision survey production and astrophysical prediction &
      Portable algorithms, verification, performance studies, and rapid
      particle--mesh method development \\
    \bottomrule
  \end{tabular}
\end{table}

\section{Cosmology model}
\label{sec:model}

\subsection{Background expansion}

The current model is a spatially flat, radiation-free $\Lambda$CDM cosmology
with one collisionless matter species.  The scale factor $a$ is normalized to
unity today and is related to redshift by $a=(1+z)^{-1}$.  With
$\Omega_\Lambda=1-\Omega_m$, the dimensionless Hubble function is
\begin{equation}
  E(a) \equiv \frac{H(a)}{H_0}
  = \left(\frac{\Omega_m}{a^3}+\Omega_\Lambda\right)^{1/2}.
  \label{eq:hubble}
\end{equation}
Radiation, curvature, massive neutrinos, evolving dark energy, and modified
gravity are not included in the present validation envelope.

Positions $\boldsymbol{x}$ are comoving and expressed in
$h^{-1}\,\mathrm{Mpc}$.  We introduce the dimensionless time
$\tau=H_0t$ and canonical momentum
\begin{equation}
  \boldsymbol{p}=a^2\frac{d\boldsymbol{x}}{d\tau}.
  \label{eq:momentum}
\end{equation}
The physical peculiar velocity is then
$\boldsymbol{v}_{\rm pec}=100\,\boldsymbol{p}/a$ in
$\mathrm{km}\,\mathrm{s}^{-1}$ for coordinates in $h^{-1}\,\mathrm{Mpc}$.

\subsection{Collisionless particle--mesh equations}

Let $\delta=\rho/\bar\rho-1$ be the matter density contrast.  We define a
rescaled potential $\phi_0$ through
\begin{equation}
  \nabla^2\phi_0 = \frac{3}{2}\Omega_m\delta,
  \qquad
  \boldsymbol{F}=-\nabla\phi_0.
  \label{eq:poisson}
\end{equation}
Using $a$ as the independent variable gives
\begin{equation}
  \frac{d\boldsymbol{p}}{da}
    = \frac{\boldsymbol{F}}{a^2E(a)},
  \qquad
  \frac{d\boldsymbol{x}}{da}
    = \frac{\boldsymbol{p}}{a^3E(a)}.
  \label{eq:eom}
\end{equation}
The periodic zero mode is explicitly removed.  Equations~\eqref{eq:poisson}
and~\eqref{eq:eom} are discretized with a uniform cell-centered mesh, equal-mass
macro-particles, and cloud-in-cell (CIC) assignment.  The same shape function is
used for scatter and gather, which avoids an asymmetric particle--mesh coupling.

One simulation step is a kick--drift--kick composition.  For a force held fixed
during a kick and a momentum held fixed during a drift, the cosmological factors
are
\begin{equation}
  K(a_0,a_1)=\int_{a_0}^{a_1}\frac{da}{a^2E(a)},
  \qquad
  D_{\rm drift}(a_0,a_1)=\int_{a_0}^{a_1}\frac{da}{a^3E(a)}.
  \label{eq:factors}
\end{equation}
Steps are uniform in $\ln a$, with the geometric midpoint separating the two
kicks.  Host-side adaptive quadrature evaluates Eq.~\eqref{eq:factors}; particle
and mesh kernels remain in the selected Kokkos execution space.  The scheme is
second order in time.  Mechanical energy is not expected to remain constant in
an expanding background; validation instead uses the linear growing mode,
convergence, mass conservation, and parallel consistency.

\subsection{Current physical limits}

The mesh supplies the only force resolution.  There is no CIC-window
deconvolution, sub-cell short-range force, adaptive mesh, or independent force
softening parameter.  The present application is therefore well matched to
linear modes and large-volume algorithm studies but not yet to precision halo
internal structure.  A credible nonlinear extension should add and separately
validate one of the following: P$^3$M, TreePM/FMM--PM, or a refined mesh.  This
choice affects conservation, scaling, load balance, and portability and should
be treated as a central research decision.

\section{Initial conditions}
\label{sec:ics}

\subsection{Linear spectrum and growth}

The $z=0$ linear matter spectrum is
\begin{equation}
  P(k)=A k^{n_s}T^2(k),
  \label{eq:pk}
\end{equation}
with $k$ in $h\,\mathrm{Mpc}^{-1}$ and $P$ in
$(h^{-1}\,\mathrm{Mpc})^3$.  The amplitude $A$ is fixed by
\begin{equation}
  \sigma_8^2=\frac{1}{2\pi^2}
  \int dk\,k^2P(k)W^2(8k),
  \quad
  W(x)=3\frac{\sin x-x\cos x}{x^3}.
  \label{eq:sigma8}
\end{equation}
The implementation supports the BBKS analytic CDM transfer function and a
CMBFAST-style table containing cold-dark-matter and baryon transfer columns.
For the table option, the total matter transfer function is formed with the
configured baryon fraction and normalized at low $k$.  The input table must
cover the largest three-dimensional mesh wave number; extrapolation beyond the
last tabulated wave number is rejected.

The growing solution is evaluated as
\begin{equation}
  D_{\rm raw}(a)=\frac{5\Omega_m}{2}E(a)
  \int_0^a\frac{da'}{a'^3E^3(a')},
  \qquad
  D(a)=\frac{D_{\rm raw}(a)}{D_{\rm raw}(1)},
  \label{eq:growth}
\end{equation}
with logarithmic growth rate $f=d\ln D/d\ln a$.  This integral solution is
cross-checked against an independent growth-ODE calculation.

\subsection{Distributed Gaussian 1LPT}

Particles begin on a cell-centered Lagrangian lattice
$\boldsymbol{q}=(\boldsymbol{i}+1/2)L/N$.  Independent Gaussian Fourier
coefficients obey
\begin{equation}
  \left\langle|\widehat\delta(\boldsymbol{k},a)|^2\right\rangle
  = \frac{D^2(a)P(k)}{L^3},
  \label{eq:fourier_variance}
\end{equation}
under the transform convention used by the implementation.  Hermitian partners
are constructed explicitly; the DC mode and initial-condition Nyquist planes
are zero.  First-order Lagrangian perturbation theory gives
\begin{align}
  \widehat{\boldsymbol{\psi}}_0(\boldsymbol{k})
    &= \frac{i\boldsymbol{k}}{k^2}\widehat\delta(\boldsymbol{k},a=1),\\
  \boldsymbol{x}(a_i)
    &= \boldsymbol{q}+D(a_i)\boldsymbol{\psi}_0,\\
  \boldsymbol{p}(a_i)
    &= a_i^2E(a_i)f(a_i)D(a_i)\boldsymbol{\psi}_0.
  \label{eq:1lpt}
\end{align}

Random values are generated from a configured 64-bit seed and a canonical
global Fourier-mode index.  Consequently, a mode receives the same coefficient
for different MPI decompositions and OpenMP schedules.  This is stronger than
merely using the same seed on every rank and enables direct phase-space
comparisons across rank counts.

The application also provides a single sinusoidal mode and a uniform lattice.
The former isolates force and growth errors; the latter tests the periodic
zero-force equilibrium.  Production cosmology will require second-order
Lagrangian perturbation theory to reduce transients and a documented path for
ingesting community initial-condition formats.

\section{A performance-portable implementation}
\label{sec:implementation}

\subsection{Reuse of the \IPPL particle--mesh stack}

The implementation follows the software decomposition established for \IPPL's
plasma particle-in-cell applications~\cite{muralikrishnan2024scaling}.  MPI
distributes fields and particles across ranks.  Kokkos views hold field and
particle attributes and Kokkos execution policies map kernels to the selected
on-node backend~\cite{trott2022kokkos}.  heFFTe supplies distributed transforms
with CPU and GPU backends~\cite{ayala2020heffte}.  Particle migration uses a
pack--send--receive--unpack path, while field boundary communication uses halo
cells.  Communication buffers can be overallocated to reduce repeated device
allocation costs.  Orthogonal recursive bisection is available when particle
clustering produces load imbalance.

This decomposition carries over directly to gravity:
\begin{enumerate}
  \item CIC scatter deposits particle mass onto the distributed mesh.
  \item The mean density is subtracted and the periodic Poisson equation is
        solved in Fourier space.
  \item The differentiated potential gives the mesh force.
  \item CIC gather interpolates the force to particle positions.
  \item Kokkos kernels apply kick and drift operations, periodic wrapping, and
        particle migration.
\end{enumerate}
No plasma-specific charge or Boris-push logic is reused.  The reusable layer is
the particle--mesh infrastructure and its parallel implementation.

\subsection{Memory, execution, and reproducibility}

Particle positions, canonical momenta, forces, global identifiers, and
Lagrangian coordinates remain in the configured Kokkos memory space during
evolution.  The background integrals are scalar host calculations.  A radial
$P(k)$ lookup is prepared on the host and copied to execution memory for initial
condition generation.  In the current prototype this lookup uses
$\mathcal{O}(N^2)$ replicated storage per rank, which is acceptable for local
validation but should be replaced or distributed before extreme-scale runs.

MPI reductions, deposition atomics, and FFT algorithms may change floating-point
summation order across decompositions and devices.  Reproducibility is therefore
defined in two layers: Fourier random phases and global particle identifiers are
deterministic; evolved floating-point fields are required to agree within
declared tolerances rather than bit for bit.  Any mixed-precision study must
report its effect on the power spectrum, force error, and growth in addition to
speed and memory savings.

\subsection{Performance methodology}

The plasma study showed that useful portability evidence separates component
costs rather than reporting only total runtime~\cite{muralikrishnan2024scaling}.
The cosmology campaign will therefore time CIC scatter, halo exchange, forward
and inverse FFTs, spectral Green-function and gradient kernels, CIC gather,
particle push, migration, load balancing, diagnostics, and output.  We will
report particle updates per second, mesh cells solved per second, time per step,
memory per particle and cell, and the fraction of runtime in communication.

Strong-scaling tests will fix particle and mesh counts; weak-scaling tests will
fix both per rank.  At least one CPU and one GPU platform should be tested, and
preferably two accelerator vendors, before claiming performance portability.
Performance-portability efficiency should be reported relative to an optimized
backend-specific baseline where available, not inferred from successful
execution alone.

\begin{table}[htbp]
  \centering
  \caption{Planned performance matrix.  Empty cells are intentional: the
  present draft does not contain accelerator or multi-node cosmology results.}
  \label{tab:performance-plan}
  \begin{tabular}{@{}p{0.25\linewidth}p{0.14\linewidth}p{0.16\linewidth}p{0.37\linewidth}@{}}
    \toprule
    Platform & Backend & Scale & Required evidence \\
    \midrule
    Local multicore CPU & OpenMP & 1--4 MPI ranks & completed correctness baseline \\
    \todo{CPU cluster} & OpenMP & multi-node & strong/weak scaling and communication split \\
    \todo{NVIDIA system} & CUDA & multi-node & kernel, FFT, migration, and end-to-end rates \\
    \todo{AMD system} & HIP & multi-node & same metrics and portability comparison \\
    \bottomrule
  \end{tabular}
\end{table}

\section{Validation}
\label{sec:validation}

\subsection{Validation hierarchy}

The validation is organized so that a failure can be attributed to one layer.
Host unit tests cover configuration, Eq.~\eqref{eq:hubble}, the growth integral
and an independent ODE oracle, Einstein--de Sitter kick and drift factors,
transfer interpolation, and $\sigma_8$ normalization.  Distributed tests then
exercise a manufactured spectral force, particle migration, and periodic
wrapping on one through four MPI ranks.

End-to-end tests include:
\begin{itemize}
  \item a uniform lattice, which must retain zero force and conserve mass and
        particle number;
  \item axis-aligned and oblique modes, which measure scale-dependent spatial
        error and linear growth;
  \item Gaussian fields, for which reconstructed Fourier coefficients test the
        requested spectrum, Hermiticity, longitudinal displacement, and the
        momentum--displacement relation;
  \item rank and thread comparisons sorted by deterministic particle ID;
  \item time-step and mesh-refinement sequences assessed independently; and
  \item matched initial-condition ensembles against an independent Zarija
        Luki\'c implementation in the common flat, Gaussian, radiation-free CDM
        subset.
\end{itemize}

The two initial-condition generators use different random-number algorithms, so
equal numeric seeds do not imply equal phases.  Their comparison therefore
separates deterministic background, transfer, normalization, and velocity-factor
checks from statistical ensemble checks.  Known lattice-origin, transform-order,
velocity-unit, and Nyquist conventions are handled explicitly rather than fitted.

\subsection{Current linear-regime results}

Table~\ref{tab:validation-results} records the current local qualification.  The
full native campaign comprised 22 simulations and 275 checks.  The matched
reference campaign comprised 84 runs and 1047 checks, including \IPPL runs on
one, two, three, and four ranks and reference runs on one, two, and four ranks.
All acceptance limits were declared before the corresponding campaign and were
not relaxed after execution.

\begin{table}[htbp]
  \centering
  \caption{Selected measurements from the local OpenMP qualification at
  initial development snapshot \texttt{0f2953bde}.  These results establish
  linear-regime behavior; later nonlinear comparisons and provenance limits
  are documented separately in Appendix~\ref{app:validation}.}
  \label{tab:validation-results}
  \begin{tabular}{@{}p{0.51\linewidth}p{0.38\linewidth}@{}}
    \toprule
    Test & Measured result \\
    \midrule
    Manufactured spectral force, ranks 1--4 & maximum relative error
      approximately $2\times10^{-15}$ \\
    Forced migration, ranks 1--4 & ownership and all IDs/attributes preserved;
      exact endpoints handled \\
    Reconstructed Gaussian Fourier coefficients & relative $L_2$ error
      $4.8\times10^{-12}$ \\
    Gaussian low-$k$ linear growth on a $32^3$ mesh & weighted complex error
      $0.333\%$ \\
    Late-$\Lambda$ single-mode growth, $z=9$ to $z=0$ & error $0.943\%$ \\
    Time convergence, position / momentum & observed order 1.988 / 1.995 \\
    Mesh refinement $16^3\rightarrow32^3\rightarrow64^3$ & growth-error
      reduction $3.899\times$ / $3.849\times$ \\
    Matched-IC ensemble mean $P_{\rm measured}/P_{\rm theory}$ & \IPPL:
      0.995736642--0.995736648; reference: 0.998164308--0.998196423 \\
    \IPPL{} rank agreement in matched campaign & maximum position RMS
      $6.502\times10^{-16}\,h^{-1}\mathrm{Mpc}$; momentum RMS
      $1.339\times10^{-17}$ in canonical units \\
    \bottomrule
  \end{tabular}
\end{table}

The result most relevant to the physical model is not any single small residual,
but the combination of second-order temporal convergence, improving spatial
error, correct Gaussian statistics, and rank-independent phases.  The remaining
sub-percent growth error at finite mesh spacing is expected for CIC PM without
window deconvolution or a short-range correction and decreases under mesh
refinement.

\subsection{Completed nonlinear comparisons and resolution controls}

The staged evidence now extends beyond this initial linear baseline.
Appendix~\ref{app:validation} documents an independent frozen-force comparison,
analytical pre-crossing pancake evolution, 32 matched nonlinear native plain-PM
runs, a separate targeted timestep extension, 34 local spatial-control runs,
and a fresh 18-run Gaussian study on Merlin CPUs.  It provides the experiment
matrices, comparison conventions, declared criteria, figures and report inventory.
The native reference uses ordinary PM evolution, not modified FastPM or COLA
stepping.

The latest Gaussian study passes every paired-code, timestep, starting-redshift,
MPI and integrity criterion, but retains 116 particle/mesh-resolution failures
among 1,355 checks.  No Gaussian shell qualifies over that full matrix.  The
local spatial study qualifies only its first low-mode shell for each fixture,
and earlier mass-sum and timestep failures remain explicitly recorded.
These results support the implementation comparison while identifying the
remaining accuracy limitations; they are not halo or continuum validation.

\subsection{Remaining scientific and performance validation}

The paper is not complete until the following evidence is added:
\begin{enumerate}
  \item Extend the present matched-PM comparison to larger crossed resolutions,
        independent realizations, converged nonlinear power spectra and halo
        statistics against an established production code.  Published
        multi-code comparisons suggest percent-level power-spectrum agreement is
        a meaningful target over a declared resolved range~\cite{angulo2022review}.
  \item Quantify force error versus separation and mesh resolution, including
        anisotropy and the effect of any future short-range solver.
  \item Add a Layzer--Irvine or equivalent expanding-universe energy diagnostic,
        while retaining mass, momentum, and particle-count checks.
  \item Validate 2LPT against an independent generator and measure transient
        suppression relative to 1LPT.
  \item Repeat correctness tests on each advertised Kokkos backend and on at
        least two multi-node systems.
  \item Populate Table~\ref{tab:performance-plan} with strong/weak scaling,
        component profiles, memory usage, and device-to-device comparisons.
  \item Replace diagnostic CSV output with scalable parallel snapshots and add
        restart, power-spectrum, light-cone or halo-catalog workflows appropriate
        to the selected science target.
\end{enumerate}

\section*{Discussion and roadmap}

The current prototype demonstrates that \IPPL's particle--mesh abstraction is
not specific to electrostatic plasma physics.  Density deposition, a global
Poisson solve, field interpolation, particle migration, and load imbalance are
also the central computational motifs of long-range cosmological gravity.  The
cosmological application can therefore reuse a portability layer whose critical
components have already been studied on heterogeneous pre-exascale systems,
while keeping cosmological equations and validation independent.

The most defensible research program has three stages.  Stage I, represented by
this draft, is a linear PM mini-app plus a staged verification suite and bounded
nonlinear comparisons, with unresolved resolution sensitivities kept explicit.
Stage II adds larger-resolution accuracy studies, 2LPT, scalable I/O and restart,
and accelerator and multi-node qualification.  Stage III is science-case dependent: a
short-range gravity method for halo resolution; neutrino or modified-gravity
particles for beyond-baseline cosmology; or coupling to a mesh hydrodynamics
solver for baryonic applications.  This staged scope retains the value of a
small portable test bed while providing a clear path toward scientifically
competitive simulations.

\section*{Code and data availability}

\IPPL is open source at \url{https://github.com/IPPL-framework/ippl}.  The
cosmology implementation and validation artifacts described here are currently
maintained on development branch \texttt{codex/cosmology-linear}; the results in
Table~\ref{tab:validation-results} correspond to commit \texttt{0f2953bde} and
its recorded development history.  Appendix~\ref{app:provenance} identifies the
later campaign reports; the Merlin run used frozen source \texttt{296fd04c0},
with documentation handoff at \texttt{9506f1e52}.  The manuscript includes
verified figure assets and their provenance manifest, not the full particle
archives.  \todo{Replace the development reference with a
public release DOI, archive input files and machine-readable results, and state
licenses before submission.}

\section*{Acknowledgements}
\todo{Funding, computing allocations, software projects, and contributors.}

\clearpage
\appendix
\numberwithin{equation}{section}
\numberwithin{figure}{section}
\numberwithin{table}{section}
\section{Validation}
\label{app:validation}

This appendix records the evidence available through 5 October 2026.  It separates
verification of equations and discrete operators, agreement between independent
implementations, and sensitivity to numerical resolution.  These are different
questions: agreement between two PM codes does not by itself establish a
continuum solution, halo accuracy, or performance portability.  All results
below concern periodic, flat, radiation-free cosmology with one collisionless
species.  Cosmological initial conditions are tested only at first Lagrangian
order; frozen-force fixtures need not be cosmological realizations.  The numerical campaigns
used an OpenMP CPU backend; no cosmological GPU, multi-node, or exascale result
is claimed.

The validation was cumulative, but it was not one homogeneous experiment.
Early campaigns were run locally on macOS; the Gaussian campaign G, broadband
campaign ZB, and subsequent GADGET-2 comparison GT are distinct executions on
the Merlin Linux login host.  Repeated cases,
reused snapshots, and related checks must not be counted as independent
statistical samples.  A \emph{check} is a recorded acceptance predicate, not an
independent measurement.  Table~\ref{tab:app-campaigns} retains the failed
predicates as well as the passing ones.  Completion of a campaign and acceptance
of all its scientific criteria are explicitly distinguished.

\begin{longtable}{@{}p{.10\textwidth}p{.35\textwidth}p{.14\textwidth}p{.31\textwidth}@{}}
\caption{Completed validation campaigns used in this appendix.  Report keys
are resolved in Section~\ref{app:provenance}.  ``Runs'' counts executable runs,
not independent realizations.  The scalar probe is a separate set of numerical
comparisons.}\label{tab:app-campaigns}\\
\toprule
Key & Experiment & Runs & Recorded outcome \\
\midrule
\endfirsthead
\toprule Key & Experiment & Runs & Recorded outcome \\
\midrule
\endhead
L & Linear evolution, rank/thread and refinement tests & 22 & 275/275 checks pass \\
ZP & Zarija background/transfer scalar probe & --- & 4,655 gated comparisons pass;
  nine additional diagnostics are not gated \\
ZI & Matched-cosmology initial conditions & 84 & 1,047/1,047 checks pass \\
F & Frozen forces versus native FastPM and independent operators & 64 &
  790/790 checks pass \\
P & Analytical pre-crossing pancake & 15 & 363/364 pass;
  one strict mass-sum check fails \\
E & Matched nonlinear native plain-PM evolution & 32 & 1,723/1,725 pass;
  two final momentum timestep checks fail \\
ER & Targeted 512-step extension of E & 2 new & 99/99 pass;
  original E failures are unchanged \\
S & Crossed particle/mesh, translation and timestep controls, local CPU & 34 &
  2,083/2,213 pass; 130 retained failures \\
G & Common-phase Gaussian controls, Merlin CPU & 18 &
  1,239/1,355 pass; 116 resolution failures \\
ZB & Broadband 1LPT evolution versus native plain-PM FastPM, Merlin CPU & 4 &
  All launches and integrity checks complete; final code differences measured, not gated \\
GT & Matched $z_i=99$ evolution with GADGET-2 TreePM, Merlin CPU & 1 integration &
  Complete from checkpoint/restart; 48 Fourier bins compared, no new acceptance gate \\
\bottomrule
\end{longtable}

\subsection{Conventions, observables, and independent checks}
\label{app:metrics}

Lengths are comoving $h^{-1}\,\mathrm{Mpc}$ and time is
$\tau=H_0t$.  Canonical momentum is $\boldsymbol p=a^2d\boldsymbol x/d\tau$;
it is neither physical peculiar velocity nor the raw velocity exported by the
Zarija generator.  Particle and mesh resolutions are denoted by $N_P$ and
$N_M$, respectively: there are $N_P^3$ particles and $N_M^3$ mesh cells, with
cell width $\Delta x=L/N_M$.  Particle comparisons are sorted by global ID,
require complete unique IDs and finite attributes, and use the minimum-image
periodic position difference.  In the matched evolution studies,
\begin{equation}
 e_x=\frac{\bigl[N_P^{-3}\sum_j|\Delta\boldsymbol x_j|^2\bigr]^{1/2}}
 {\Delta x},\qquad
 e_p=\frac{\bigl[\sum_j|\Delta\boldsymbol p_j|^2\bigr]^{1/2}}
 {\bigl[\sum_j|\boldsymbol p_{j,\mathrm{ref}}|^2\bigr]^{1/2}}.
 \label{eq:app-phase}
\end{equation}
The analytical pancake instead normalizes trajectory errors by the corresponding
analytical displacement or momentum RMS; its percentages should not be confused
with position errors in mesh cells.

For the direct-mode nonlinear comparisons E, S, and G, nonzero density modes
are measured directly from the particle positions, independently of either
force mesh:
\begin{equation}
 \widehat\delta_{\boldsymbol n}=
 \frac{1}{N_P^3}\sum_j\exp\!\left[-\frac{2\pi i}{L}
 \boldsymbol n\!\cdot\!\boldsymbol x_j\right].
 \label{eq:app-direct}
\end{equation}
For the vectors $\boldsymbol d_A,\boldsymbol d_B$ of coefficients in a common
mode set, define $S_A=\|\boldsymbol d_A\|_2^2$ and $S_B=\|\boldsymbol d_B\|_2^2$.
The reported power discrepancy, symmetric complex residual, and signed
cross-correlation are
\begin{equation}
 e_P=\left|\frac{S_A}{S_B}-1\right|,\qquad
 e_\delta=\frac{\|\boldsymbol d_A-\boldsymbol d_B\|_2}{(S_AS_B)^{1/4}},
 \qquad
 r=\frac{\operatorname{Re}(\boldsymbol d_B^\dagger\boldsymbol d_A)}
 {\sqrt{S_AS_B}}.
 \label{eq:app-density}
\end{equation}
Absolute differences are retained when the signal is undefined; relative
metrics are not manufactured by flooring an absent signal.  In these density
comparisons the saved definition treats $\min(S_A,S_B)\le10^{-28}$ as undefined.
For plotted shell powers in E, S, and G,
$P(k)=L^3\langle|\widehat\delta|^2\rangle$ and no shot-noise subtraction is
applied.  The separate ZB shell estimator uses periodic CIC assignment, applies
the CIC window correction, and subtracts Poisson shot noise as described in
Section~\ref{app:zeldovich-fastpm}.

Independent tests exercise the background growth ODE, Einstein--de Sitter
integrals, transfer interpolation, normalization quadrature, manufactured
spectral forces, CIC scatter/gather, periodic wrapping, and particle migration.
The wrapping regression includes exact box endpoints and excursions by multiple
box lengths; ownership, IDs and migrated attributes are checked on ranks 1--4.
Manufactured spectral-force errors are approximately $2\times10^{-15}$.
These component tests isolate implementation defects before a dynamical
comparison is interpreted as evidence about cosmology.

\subsection{Linear evolution and refinement}
\label{app:linear}

Campaign L uses $L=200\,h^{-1}\mathrm{Mpc}$,
$(\Omega_m,\Omega_b,h,n_s)=(0.31,0.0487,0.675,0.965)$ and BBKS transfer.
The Gaussian case deliberately sets $\sigma_8=0.02$ to remain linear; it is
not the later $\sigma_8=0.82$ nonlinear experiment.  Most runs evolve from
$z=49$ to $24$.  Axis-aligned and oblique $(1,1,1)$ perturbations start with
density amplitude $10^{-3}$; the late-$\Lambda$ test evolves $z=9$ to $0$
with amplitude $10^{-4}$.  Uniform, axial and Gaussian cases cover ranks 1--4,
and a separate one-rank/two-thread run checks OpenMP reproducibility.

All uniform configurations retain exactly zero displacement, momentum and
force.  Axis/oblique growth errors are approximately $0.1561\%/0.1578\%$;
late-$\Lambda$ growth differs from the independent reference by $0.9427\%$.
The Gaussian low-$k$ complex-growth error over 13 selected modes is $0.3332\%$,
and independent reconstruction of the initial Gaussian coefficients gives
relative $L_2$ error $4.804\times10^{-12}$.  The largest recorded mass residual
is $1.686\times10^{-12}$, below the declared $2\times10^{-12}$ criterion.

The 4/8/16-step sequence measures successive-solution temporal orders
1.9881 in position and 1.9954 in momentum.  Joint particle/mesh refinement
$N_P=N_M=16,32,64$ at 64 steps reduces the growth errors from
$5.9551\times10^{-3}$ to $1.5275\times10^{-3}$ and $3.9681\times10^{-4}$,
with reduction factors 3.8985 and 3.8495.  This is evidence for the tested
linear discretization, not an independent fixed-particle mesh study.
The early report contains the measured outcomes but lacks the executable/source
hash linkage added in later campaigns; a subsequent rebuilt executable must not
be described as cryptographically identified by that report.

\subsection{Matched initial conditions against Zarija}
\label{app:zarija}

\paragraph{Matched physical subset and deterministic probes.}
Campaigns ZP and ZI compare the original Zarija Luki\'c generator with \IPPL
only in their common flat, Gaussian, radiation-free CDM/1LPT subset:
$\Omega_\nu=\Omega_r=f_{\mathrm{NL}}=0$ and $w=-1$.  BBKS and a supplied
CMBFAST-style table correspond to reference transfer flags 4 and 0.
The original reference sources and transfer input were preserved.
The scalar probe evaluates the public reference cosmology implementation,
independently of particle statistics.  Its 4,655 gated comparisons pass;
the additional nine excluded DC/legacy-table diagnostics are not passing
acceptance tests (one BBKS DC diagnostic and eight table-interval diagnostics).

Across the sampled simulation wavenumbers, transfer discrepancies are below
$5\times10^{-16}$, against a $2\times10^{-12}$ gate.  Power-normalization
differences are $4.763\times10^{-7}$ for BBKS and $3.255\times10^{-5}$ for the
table, against $5\times10^{-4}$.  At $z=0,9,49,200$, maximum relative differences
in $D$, $dD/d\tau$ and $f$ are $3.334\times10^{-7}$,
$7.255\times10^{-7}$ and $3.921\times10^{-7}$, respectively, below
$2\times10^{-5}$.  The reference's finite-start growth solution and numerical
integration accuracy are distinct from the \IPPL growing-mode quadrature.

\paragraph{Ensemble design and convention conversion.}
The particle experiment fixes $N_P=N_M=32$, $L=168.75\,h^{-1}\mathrm{Mpc}$ and
$(\Omega_m,\Omega_b,h,\sigma_8,n_s)=(0.31,0.0487,0.675,0.82,0.965)$.
Two transfers, starting redshifts 49 and 200, and eight seeds are used:
1234567, 7654321, 104729, 130363, 32452843, 49979687, 67867967 and 86028121.
The 64 serial realizations are augmented by 12 \IPPL runs on ranks 2, 3 and 4,
and eight reference runs on ranks 2 and 4, using the first seed for each
transfer/redshift combination.  Thus the 84 runs comprise 44 \IPPL and 40
reference executions.  The uneven three-rank case is an \IPPL test, not a
claimed reference slab-decomposition test.

Different generators use different random-number algorithms, so identical
numeric seeds do not imply equal cross-code phases.  Cross-code agreement is
therefore assessed statistically.  Within-code changes in redshift and transfer
use phase-matched comparisons.  Rank changes are phase matched only for \IPPL;
the reference's MPI runs test statistics because its RNG realization changes
with decomposition.  ID-derived Lagrangian
coordinates, the node/cell-centred lattice offset, the half-cell Fourier phase,
transform ordering and velocity units are handled explicitly.  For the reference
velocity $\boldsymbol v_Z$, the conversion is
$\boldsymbol p=a^2\boldsymbol v_Z/100$; the physical peculiar velocity is
$100\boldsymbol p/a=a\boldsymbol v_Z$ in $\mathrm{km}\,\mathrm{s}^{-1}$.

The analysis reconstructs displacement and then the $z=0$ density coefficients
from $-i\boldsymbol k\cdot\widehat{\boldsymbol\psi}_0$; it does not infer the
linear spectrum from nonlinear Eulerian CIC density.  It excludes DC and every
componentwise Nyquist plane and retains one member of each conjugate pair:
14,895 independent complex modes per seed, or 119,160 per eight-seed ensemble.
Redshift-reused realizations are not additional independent samples.

\paragraph{Results and statistical meaning.}
All 1,047 checks pass.  Ensemble means of modewise recovered/theoretical power
are approximately 0.995737 for \IPPL and 0.99816--0.99820 for the reference.
Six-standard-error criteria are based on the expected Gaussian statistics;
the reference additionally carries a predeclared $5\times10^{-4}$ deterministic
normalization allowance.  The mean-power half-widths are 0.01738145 and
0.01788145, respectively.  Thus these means are consistent with the expected
finite-ensemble scatter; they are not a fitted amplitude correction.
Figure~\ref{fig:app-ic-power} distinguishes theoretical one-standard-error bars
from the wider acceptance band.

Maximum canonical-momentum relation errors are $1.998\times10^{-12}$ for
\IPPL and $3.921\times10^{-7}$ for Zarija.  Longitudinal-displacement residuals
remain below $2.111\times10^{-13}$ in both.  Same-transfer redshift-scaling
errors are $6.263\times10^{-13}$ and $5.138\times10^{-8}$; transfer-scaling
errors are $3.169\times10^{-9}$ and $1.604\times10^{-5}$.
The 12 \IPPL MPI comparisons give maximum componentwise position RMS
$6.501\times10^{-16}\,h^{-1}\mathrm{Mpc}$ and componentwise momentum RMS
$1.339\times10^{-17}$ in canonical units
(Figure~\ref{fig:app-ic-mpi}).

\paragraph{An explicit legacy exclusion.}
The original table's first weighted transfer entry, about $1.9965\times10^7$,
is left unnormalized by the reference.  Independent integration of this narrow,
malformed low-$k$ interval contributes 1.6618 times the legacy routine's reported
total raw variance, revealing that its quadrature misses the discontinuity.
Agreement on the simulation band therefore does not validate that interval.
A separately normalized copy diagnoses the issue without modifying the original
input.  Different RNGs, Nyquist handling, this interval, and untested reference
physics are exclusions from equivalence; the experiment does not establish
feature parity with the whole Zarija code or identical production input formats.

\validationplot{ic-power.png}{Recovered $z=0$ Lagrangian power from eight
independent realizations per code at $z_i=49$, for BBKS and tabulated transfers.
Lower panels show ensemble means of modewise power ratios, not ratios of shell
means.  Error bars are theoretical Gaussian one-standard-error uncertainties;
the shaded band is the saved six-standard-error reference criterion including
its normalization allowance.  DC and componentwise Nyquist planes are excluded;
the highest shells have incomplete angular coverage.  Source: ZI/ZP.}
{fig:app-ic-power}

\validationplot{ic-mpi.png}{ID-matched \IPPL initial-position and canonical-
momentum componentwise RMS differences on ranks 2--4 relative to one rank, normalized by the
declared limits.  Both transfers and initialization redshifts are included.
This is a decomposition-reproducibility test, not a strong-scaling, GPU or
throughput measurement.  Source: ZI.}{fig:app-ic-mpi}

\subsection{Frozen-force comparison with native FastPM}
\label{app:force}

The frozen experiment isolates scatter, the mesh force, and gather without
advancing particles.  Reference FastPM~\cite{feng2016fastpm} is pinned at commit
\texttt{15b6c4fd7502a81d99dd13f54fcc9cfa44be1331}; the tracked upstream sources
are unchanged.  Eight fixtures on ranks 1--4 in both codes give 64 runs.
They include a uniform state, a strictly common-band mode, axis and oblique
deformations, subcell shifts, shuffled/wrapped jitter, a dense cluster, a second
mesh size, and an Einstein--de Sitter case.  Most use $N=16$,
$L=168.75\,h^{-1}\mathrm{Mpc}$ and $\Omega_m=0.31$; the matrix also contains
$N=32$ and $\Omega_m=1$ cases.

The compared force is that of Eq.~\eqref{eq:poisson}, with CIC scatter and
gather and no window deconvolution.  Native FastPM uses its \texttt{NAIVE}
spectral kernel without softening.  The node-mesh origin receives the explicit
shift $\boldsymbol x_F=\operatorname{wrap}(\boldsymbol x_I-L/(2N_M))$ in each
coordinate.  Its exported native acceleration is multiplied by
$-3\Omega_m/2$ to obtain the common force convention; this is a documented
unit/sign conversion, not a fitted normalization.

The full native discrete operators are \emph{not identical}.  \IPPL zeros the
differentiated component's Nyquist plane, whereas native FastPM zeros only the
eight self-conjugate corners and retains its native real-to-complex behavior.
FastPM FFT, mesh and CIC operations are double precision, but its wavevector
tables and particle accelerations use float32.  Independent NumPy constructions
of CIC deposition, each complete native spectral operator and reciprocal gather
test each code against its own oracle.  The expected difference between these
operators is then compared with the actual residual field.  A diagnostic
common-band projection is never inserted into either production solve.

All 790 checks pass.  Maximum relative RMS particle-force errors against the
respective independent oracles are $6.741\times10^{-14}$ for \IPPL and
$3.205\times10^{-8}$ for FastPM; mesh-force errors are
$5.672\times10^{-14}$ and $6.474\times10^{-14}$.
The specially constructed common-band fixture has raw particle-force agreement
within $4.142\times10^{-8}$ over ranks 1--4.  In contrast, raw oblique and
wrapped-jitter differences reach $2.467\%$ and $11.714\%$ and are predicted by
the preserved Nyquist operators.  These are not hidden by a looser raw-force
gate.  Uniform relative-force errors are undefined when both forces are zero.
Figure~\ref{fig:app-force} reports residual fields, not a difference between
the separate force magnitudes.

\validationplot{frozen-force.png}{Frozen operators on ranks 1--4.  Each native
code is compared with its own independent CIC/spectral/gather oracle; raw and
predicted cross-code differences retain the different Nyquist conventions.
Values are maxima over tested ranks.  Exactly zero uniform-force cases are
omitted from relative logarithmic plots, not assigned a positive floor.
Passing this test does not make arbitrary unfiltered native forces identical.
Source: F.}{fig:app-force}

\subsection{Analytical pancake before shell crossing}
\label{app:pancake}

An independent planar growing-mode solution tests actual evolution, not just
the force at a prescribed state.  Before crossing,
\begin{equation}
 \boldsymbol x=\boldsymbol q-
 A(a)\sin(\boldsymbol k\cdot\boldsymbol q)\frac{\boldsymbol k}{k^2},
 \qquad
 \boldsymbol p=a^2E(a)f(a)(\boldsymbol x-\boldsymbol q),
 \label{eq:app-pancake}
\end{equation}
where $A(a)$ follows independently evaluated growth.  The Jacobian along the
normal is $J=1-A\cos(\boldsymbol k\cdot\boldsymbol q)$, and
$\rho/\bar\rho=J^{-1}$ while the map is single valued.  Final amplitudes 0.5
and 0.8 give peak density contrasts 1 and 4; the analytic oracle is never
continued past crossing.

Campaign P evolves $z=49$ to $9$ and varies coupled particle/mesh sizes
16/32/64, timestep counts 16/32/64/128, axis/oblique orientation and ranks 1--4.
Spatial errors compare with continuum trajectories at 128 steps; temporal
errors compare successive numerical solutions at fixed mesh.  At $N=64$ and
$A_f=0.8$, displacement and momentum relative RMS errors are
$0.413790\%$ and $0.735828\%$.  Spatial orders span 1.4337--1.7603 and temporal
orders 1.7994--2.0183.  These results pass their engineering trajectory budgets
but do not establish second-order spatial accuracy
(Figure~\ref{fig:app-pancake-convergence}).

One of 364 checks fails: the mass diagnostic reaches
$2.07200923\times10^{-12}$ against $2\times10^{-12}$ at step 89,
$a=0.0612396326$, for expected total unit mass 262,144.  The absolute reported
discrepancy is $5.43165\times10^{-7}$.  Independent endpoint CIC redepositions
summed with \texttt{math.fsum} recover exactly 262,144, but the worst intermediate
state was not saved.  Summation roundoff is therefore a plausible explanation,
not a proved resolution of that failed gate.  Neither a summation-error bound
nor the user's acceptance of a practical engineering baseline changes the
original failed status.

Local gradients are a separate limitation.  At $A_f=0.8$, sampled minimum
Jacobians change $0.207978\to0.181079\to0.172239$ on 16/32/64, rather than
converging monotonically to the continuum value 0.2.  Maximum errors against
the identically finite-differenced analytical map remain about 0.19.
Figure~\ref{fig:app-pancake-local} shows why shrinking support of a trajectory
defect can improve a global RMS while leaving local density unresolved.
Grid/lattice locking is a possible explanation, not an established cause.

\validationplot{pancake-convergence.png}{Analytical pre-crossing pancake:
global continuum trajectory errors under joint particle/mesh refinement and
successive-solution differences under fixed-mesh temporal refinement.  Position
and momentum use their respective analytical RMS scales.  Slope guides are
not fitted results.  The original mass-sum failure is retained; passing global
trajectory tests does not establish local-density accuracy.  Source: P.}
{fig:app-pancake-convergence}

\validationplot{pancake-local.png}{Local displacement and interval-Jacobian
errors along one transverse Lagrangian row.  Numerical and exact maps are
finite-differenced identically.  These are not Eulerian CIC density curves.
The persistent local defects explain why the global trajectory results alone
cannot qualify peak density or caustic structure.  Source: P and its saved
diagnostic audit.}{fig:app-pancake-local}

\subsection{Matched nonlinear evolution and the targeted time refinement}
\label{app:evolution}

Campaign E advances identical imported positions and canonical momenta with
the production \IPPL force/KDK methods and native
\texttt{fastpm\_solver\_evolve} in \texttt{FASTPM\_FORCE\_PM} mode.
It does not use the modified FastPM stepping scheme or COLA.  Initial momenta
are rounded once to float32 and those same values are supplied to both codes;
native momentum/acceleration remain float32 and positions/FFTs remain double.
Origin shifts are reversed when actual native states are exported.
Logarithmic full-step endpoints, geometric kick midpoints and synchronized
native outputs are compared without snapshot interpolation.  Actual native
kick/drift factors agree with independent quadrature within
$6.740\times10^{-14}$, below the $10^{-8}$ gate.

The 32-run matrix fixes $N_P=32$, $L=168.75\,h^{-1}\mathrm{Mpc}$,
$\Omega_m=0.31$ and $a=0.02\to0.2$.  A planar fixture has extrapolated final
linear amplitude 1.5 and must exhibit actual sheet folding; a second fixture
contains nine coupled three-dimensional growing modes with final linear
density RMS 1.  The latter is a deterministic nonlinear test, not a
$\sigma_8$-normalized Gaussian realization.  Meshes 16/32/64, 64/128/256 steps
and ranks 1--4 are tested, with nine synchronized checkpoints per run.
Unequal particle/mesh counts use the explicit cell-mass normalization
$N_M^3/N_P^3$ before subtracting unity.

The accepted density band consists of 128 unique conjugate pairs with
$0<|\boldsymbol n|\le4$, or four active harmonics for the planar fixture,
measured with Eq.~\eqref{eq:app-direct}.  Aggregate power/complex budgets are
5/2/1\% for $N_M=16/32/64$, with minimum correlations 0.995/0.999/0.9995.
Planar cross-code position and momentum budgets are $10^{-3}$ cells and
$10^{-3}$ relative RMS.  Raw three-dimensional trajectories and individual
shells are characterized rather than gated as identical-operator solutions.

All imported-state, paired-code, rank, native-factor, mean-momentum and folding
checks pass.  Planar cross-code differences stay below
$1.719\times10^{-6}$ cells and $3.821\times10^{-7}$ relative momentum.
Coupled-3D aggregate power/complex differences reach only
$0.24245\%/0.33780\%$, with $r\ge0.999994495$.
Raw 3D particle differences nevertheless reach 0.03205 cells and $1.0072\%$
relative momentum.  Figures~\ref{fig:app-phase-space} and
\ref{fig:app-nonlinear-modes} illustrate this matched discrete comparison,
not an analytical solution after crossing.

The two failures in E are final 128-to-256-step momentum differences:
$0.276343\%$ for \IPPL and $0.275555\%$ for native FastPM, exceeding the
unchanged $0.2\%$ criterion despite reduction factors near four.  The separate
two-run ER extension adds 512 steps only at $N_P=N_M=32$ on one rank.
Its 256-to-512 differences are $0.0696213\%/0.0694281\%$ in momentum and
$0.000449258/0.000448378$ cells in position; all 99 checks pass, with final
observed orders 1.9725--2.0083 for the assessed 3D quantities.  Below the
predeclared analysis floors no order is inferred.  ER does not relabel E as
all-pass and supplies no unperformed 512-step MPI comparison.

Fixed-particle mesh sensitivity is not continuum convergence: the pancake
power changes increase from $0.5367\%$ on 16-to-32 meshes to $3.6296\%$ on
32-to-64.  The maximum \IPPL CIC mass residual in E is
$1.219\times10^{-12}$; that smaller residual does not overturn P's distinct
failed mass check.  Native unit-particle counts and \IPPL deposited-CIC sums
are different diagnostics and are never treated as interchangeable measures.

\validationplot{nonlinear-phase-space.png}{Matched particle-sheet phase
portraits at three saved epochs, including the folded final sheet.  Curves
connect the same 32 ID-matched particles in one transverse Lagrangian row,
with breaks at periodic edges; they are not an interpolated distribution.
The pre-crossing analytical result is not an oracle for the final folded state.
Source: original 32-run campaign E.}{fig:app-phase-space}

\validationplot{nonlinear-modes.png}{Resolved-density evolution from exact
particle sums, with $N_P=32$, meshes 16/32/64, 256 steps and one rank.  The
3D comparison uses 128 unique low-mode pairs; the pancake uses four
active harmonics.  The figure retains E's two failed timestep checks, which
are not failures of the plotted cross-code density comparison.  The separate
512-step follow-up ER is described in the text, not substituted into this
figure.}{fig:app-nonlinear-modes}

\subsection{Crossed spatial resolution and translation controls}
\label{app:spatial}

Campaign S removes an ambiguity of earlier tests by varying particle sampling
and force-mesh size independently.  For each planar/coupled-3D fixture, it
crosses $N_P=32,64$ with $N_M=32,64$ at 1,024 steps, adds rigid translations
at $N_P=64$ on both meshes, and adds 512/2,048-step cases at $N_P=N_M=64$.
The coupled-3D finest standard case is also repeated on three ranks.  Seventeen
paired cases give 34 runs, with the same $L$, $\Omega_m$ and $a=0.02\to0.2$
as E and nine saved epochs per run.  The fixed physical translation is
$\boldsymbol\Delta=(0.37,0.23,0.41)L/64$; its known Fourier phase is removed
before comparing coefficients.

Direct modes with $0<|\boldsymbol n|\le4$ are separated into three shells:
$0<|\boldsymbol n|<1.5$, $1.5\le|\boldsymbol n|<2.5$, and
$2.5\le|\boldsymbol n|\le4$.  A bookkeeping edge at 4.5 does not extend the
measured band beyond 4.  For campaign G only, fixed $128^3$ interlaced PCS
extraction, with its assignment-window correction, characterizes higher modes
through 12.  Its
low-mode coefficients must independently agree with the direct sums within
$\max(10^{-12},10^{-3}\|\boldsymbol d_{\rm direct}\|_2)$.
PCS is an analysis tool; the simulation forces remain CIC.

\begin{table}[htbp]
\centering
\caption{Predeclared shell-wise controls in S and G.  Percentages denote
$e_P$ or $e_\delta$ in Eq.~\eqref{eq:app-density}, not errors against truth.
All applicable saved epochs and both codes enter the controls.}
\label{tab:app-budgets}
\begin{tabular}{@{}lrrr@{}}
\toprule Control & $e_P$ limit & $e_\delta$ limit & Minimum $r$ \\
\midrule
Paired codes, $N_M=32$ & $2\%$ & $2\%$ & 0.999 \\
Paired codes, $N_M=64$ & $1\%$ & $1\%$ & 0.9995 \\
Particle or mesh refinement & $5\%$ & $5\%$ & 0.999 \\
Dephased translation (S) & $1\%$ & $2\%$ & 0.9998 \\
Starting redshift (G, final epoch) & $2\%$ & $3\%$ & 0.999 \\
Finest timestep shell difference & --- & $0.2\%$ & --- \\
\bottomrule
\end{tabular}
\end{table}

Qualification requires a contiguous prefix of shells for which \emph{every}
applicable control passes.  Failed integrity, extraction, rank or global-time
controls veto qualification.  Passing an aggregate low-mode norm or choosing a
favorable final epoch is not an alternative to this rule.

Of 2,213 checks, 130 fail: 56 particle-resolution, 51 mesh-resolution,
20 translation and three paired-code shell checks.  Both fixtures qualify only
the first shell, $0<|\boldsymbol n|<1.5$, over the declared matrix.
For the 3D fixture the largest measured mode in that shell is $\sqrt2$;
for the planar fixture it is the fundamental mode.  All required global,
timestep and rank controls pass.

The maximum power sensitivity is $16.77\%$ for the pancake and $19.27\%$ for
coupled 3D.  The finest-mesh translated pancake reaches $1.336\%$ power
difference versus $1\%$; coupled-3D translation reaches $4.160\%$ complex
difference versus $2\%$.  The three cross-code failures occur only for the
undersampled $N_P=32,N_M=64$ coupled-3D third shell at intermediate epochs:
complex differences of approximately 1.11--1.39\% exceed 1\%, while the power
and correlation conditions pass.  Close matched $64^3/64^3$ agreement does
not remove these failures from the crossed matrix.

Final 1,024-to-2,048-step coupled-3D momentum differences are approximately
$2.683\times10^{-4}$ relative, with near-second-order late-time reduction.
Some pancake differences are below the predeclared analysis floor, and no
order is inferred for them.  Three-rank differences stay below
$5.9\times10^{-15}$ cells/$5.1\times10^{-15}$ relative momentum for \IPPL
and $1.16\times10^{-7}$ cells/$8.63\times10^{-8}$ for native PM.
Four separate \IPPL pancake runs also exceed the historical $2\times10^{-12}$
mass diagnostic, with maximum $2.951\times10^{-12}$.  These explicitly retained
historical flags are not recast as passing tests or added to the new shell-gate
count.  Figure~\ref{fig:app-spatial} shows why code agreement alone would
have missed the particle-sampling and grid-alignment sensitivities.

\validationplot{spatial-controls.png}{Independent particle and force-mesh
controls on the local CPU platform.  Rows show planar and coupled-3D fixtures.
Power/translation panels are worst-case envelopes over both codes and nine
saved epochs, not ensemble uncertainties or continuum errors.  Translation
phases are removed analytically.  Temporal panels compare 512/1,024/2,048 steps
at $N_P=N_M=64$.  Only the first exact shell meets every declared control;
all 130 failures remain recorded.  Source: S.}{fig:app-spatial}

\subsection{Common-phase Gaussian evolution on Merlin}
\label{app:gaussian}

Campaign G is a fresh Linux/OpenMP experiment, not a resumed macOS journal.
It uses one mode-keyed Gaussian BBKS realization (seed 20261003) with common
physical phases at both particle resolutions, initial band
$0<|\boldsymbol n|\le12$, $L=168.75\,h^{-1}\mathrm{Mpc}$, and
$(\Omega_m,\Omega_b,h,n_s,\sigma_8)=(0.31,0.0487,0.675,0.965,0.82)$.
The BBKS choice has no baryonic acoustic transfer features; the baryon parameter
does not make this a full baryon-physics experiment.  $\sigma_8$ normalizes the
continuous spectrum, not the measured finite box or truncated realization.
Both codes receive the identical imported, rounded-once canonical momenta.

Nine paired cases comprise the four crossed $N_P,N_M=32,64$ combinations
at $z_i=49$, 2,048 steps; finest-grid 1,024/4,096-step controls at $z_i=49$;
finest-grid 2,048/4,096-step controls at $z_i=99$; and the standard finest case
on four ranks.  All 18 runs end at $a=1$, saving nine synchronized epochs.
Starting-redshift comparisons use only the final epoch at 4,096 steps:
intermediate indices denote different scale factors for different starts.
No interpolation is used.

The build uses GCC 14.3, OpenMPI 5.0.10, Kokkos 5.2.0, heFFTe 2.4.1 and
the pinned native reference.  Metadata records OpenMP/Host execution and
actual MPI size: at most four ranks/workers, one thread per nonlinear rank,
no GPUs.  Login-host timings are not performance results.  All 21 cosmology
CTests pass, including ranks 1--4; an eight-run pipeline smoke passes 364 checks.
Neither substitutes for the full scientific matrix.

All 1,355 checks are present: 1,239 pass and 116 fail, equally divided between
particle and mesh resolution.  Every integrity, measurement, paired-code,
timestep, start-redshift and rank check passes, but \emph{none of the three
shells qualifies across the complete resolution matrix}.  Lowest-shell power
sensitivities reach $6.246\%$ (particles) and $8.137\%$ (mesh) at intermediate
epochs, exceeding 5\%; its complex and correlation criteria pass.
Higher-shell power changes reach $27.91\%$ and $30.61\%$, respectively, in
both implementations.  No retrospective subset or relaxed tolerance removes them.

\begin{table}[htbp]
\centering
\caption{Gaussian CPU results in G.  Paired-code values concern the exact
direct band $0<|\boldsymbol n|\le4$.  Finest-case values are at $a=1$,
$N_P=N_M=64$, $z_i=49$, 2,048 steps.  Temporal/rank differences are not truth
errors.}
\label{tab:app-gaussian}
\small
\begin{tabular}{@{}>{\raggedright\arraybackslash}p{.55\linewidth}>{\raggedright\arraybackslash}p{.40\linewidth}@{}}
\toprule Observable & Result \\
\midrule
Largest paired-code power / complex difference & $0.211604\%$ / $0.802740\%$ \\
Minimum paired-code correlation & 0.999967822 \\
Finest-case power / complex difference & $0.00120759\%$ / $0.00362866\%$ \\
Finest-case minimum correlation & 0.999999999342 \\
Starting-$z$ power / complex sensitivity & $1.463041\%$ / $2.649273\%$ \\
Largest final 2,048-to-4,096 position difference & $3.970\times10^{-6}$ cells \\
Largest final 2,048-to-4,096 momentum difference & $2.448\times10^{-6}$ relative \\
Largest shell temporal complex difference & $1.303\times10^{-6}$ \\
Rank 1 versus 4, \IPPL position / momentum & $3.14\times10^{-15}$ cells /
  $1.26\times10^{-15}$ relative \\
Rank 1 versus 4, native PM position / momentum & $2.33\times10^{-8}$ cells /
  $1.85\times10^{-8}$ relative \\
\bottomrule
\end{tabular}
\end{table}

Final $z_i=49$ position orders are 2.0003 for \IPPL and 1.7405 for native PM.
The corresponding final momentum/density differences are below the predeclared
analysis floor; this does not demonstrate machine-precision saturation.
The two $z_i=99$ step counts measure a refinement difference, not an order.
The largest Gaussian \IPPL mass residual is $5.018\times10^{-14}$, without
changing any earlier mass failure.  The practical conclusion of
Figure~\ref{fig:app-gaussian} is strong matched-discretization agreement but
insufficient particle/mesh robustness in the tested 32/64 matrix.

\validationplot{gaussian-controls.png}{Merlin OpenMP evolution of one common
band-limited Gaussian BBKS realization.  Top: final spectra and paired-code
residuals for the crossed 32/64 grids at $z_i=49$, 2,048 steps.  Filled symbols
use exact modes $|\boldsymbol n|\le4$; open symbols and the higher shaded band
are PCS characterization only.  Bottom: worst-epoch spatial sensitivities,
starting-redshift changes at 4,096 steps, and time refinement.  All 116
resolution failures remain visible, leaving zero qualified shells.  Source: G.}
{fig:app-gaussian}

\subsection{Broadband Zel'dovich evolution against native FastPM}
\label{app:zeldovich-fastpm}

Campaign ZB follows the low-band frozen-force and nonlinear checks with a
matched, broadband particle evolution.  One Gaussian realization (seed
20261003) supplies common 1LPT initial conditions at $z_i=99$ and $49$.
The spherical Fourier band is $0<|\boldsymbol n|\le48$, or 75\% of the
particle Nyquist frequency for $N_P=128$.  The periodic box is
$L=168.75\,h^{-1}\,\mathrm{Mpc}$ with a $128^3$ particle load and a $128^3$
CIC force mesh.  The recorded background uses $(\Omega_m,h,n_s,\sigma_8)=
(0.31,0.675,0.965,0.82)$ and a pure BBKS transfer.  Although a baryon density
is recorded in the input metadata, BBKS here has no baryon-acoustic features;
the experiment does not model radiation, neutrinos, gas, or baryonic effects.

Both solvers import the same once-rounded float32 particle momenta, without
growth rescaling.  \IPPL uses the production particle--mesh kick--drift--kick
adapter; the pinned native FastPM run uses ordinary plain-PM evolution, not
FastPM's modified fast-force stepping or COLA.  Thus ZB is a cross-code
comparison of a closely matched PM model, not a comparison with GADGET-2 or an
independent short-range gravity method.  Each case advances to $a=1$ with 24
uniform-in-$\ln a$ intervals (2,400 steps from $z_i=99$ and 2,040 from
$z_i=49$); the two starts have nearly equal step size in $\ln a$.  The change
in start epoch and step count is a sensitivity comparison, not a convergence
order test.

Four launches ran sequentially on the Merlin login CPU using eight
single-threaded MPI ranks.  All completed under the fixed 1,200-second
per-launch limit.  Each of the four runs saved 25 checkpoints on eight ranks;
all 800 archived rank shards, output metadata, source/executable provenance,
and input hashes verified.  The imported initial positions round-trip within
$7.3\times10^{-16}$ force-mesh cells and momenta match exactly.  The largest
recorded \IPPL checkpoint mass residual is $7.11\times10^{-15}$.  These are
integrity and operational checks; no threshold was imposed on the final
cross-code scientific differences.

At each saved epoch, density was assigned periodically by CIC to the $128^3$
analysis mesh, the CIC power window was corrected, and Poisson shot noise was
subtracted.  Shells include modes through $|\boldsymbol n|=48$.  Table
\ref{tab:app-zeldovich-fastpm} reports final particle-phase-space differences
and the largest absolute binwise $P(k)$ ratio difference among the 48 shells
where both shot-noise-subtracted powers are positive.

\begin{table}[htbp]
\centering
\small
\caption{Final $z=0$ cross-code comparison in ZB.  Momentum differences are
relative to native FastPM.  The power maximum is over 48 positive
shot-noise-subtracted shells through the initial-condition mode ceiling.}
\label{tab:app-zeldovich-fastpm}
\begin{tabular}{@{}rrrrr@{}}
\toprule
$z_i$ & Position RMS [$\Delta x$] & Momentum RMS [relative] &
  $\max|P_{\mathrm{IPPL}}/P_{\mathrm{FastPM}}-1|$ & Positive shells \\
\midrule
99 & $6.0085\times10^{-3}$ & $1.7052\times10^{-3}$ & $6.6987\times10^{-5}$ & 48 \\
49 & $5.1486\times10^{-3}$ & $1.6049\times10^{-3}$ & $5.7355\times10^{-5}$ & 48 \\
\bottomrule
\end{tabular}
\end{table}

The maximum power differences are 0.00670\% ($z_i=99$) and 0.00574\%
($z_i=49$).  This close agreement is consistent with the shared realization
and closely matched PM force and integration path.  It does not establish
agreement with a continuum solution: the run has one realization, uses a
single particle/mesh resolution, and its high-$k$ bins are characterization
near the particle Nyquist, not a qualified accuracy band.  The logarithmic
power spectrum and ratio are shown in Figure~\ref{fig:app-zeldovich-fastpm}.
The left panel is the older cutoff-12 FastPM-only resolution study and is kept
separate from the new matched broadband measurements.

\validationplot{zeldovich-fastpm-comparison.png}{Matched broadband 1LPT
evolution at $z=0$.  Center: periodic-CIC, CIC-window-corrected,
shot-noise-subtracted spectra from the shared-realization $128^3$ runs and the
linear BBKS spectrum.  Right: binwise \IPPL/native plain-PM FastPM power ratio
difference for the two starting redshifts.  The left panel preserves the
historical cutoff-12 FastPM-only resolution study; it is not merged with ZB.
The high-$k$ points are characterization, not a convergence claim.  Source:
ZB.}{fig:app-zeldovich-fastpm}

\subsection{Matched initial-redshift-99 evolution with GADGET-2 TreePM}
\label{app:gadget2-zeldovich}

Campaign GT adds an independently implemented gravity solver to the saved
ZB $z_i=99$ comparison.  It is a matched-realization comparison, not a
random-seed-matched ensemble: IPPL and native FastPM use their common saved
Gaussian 1LPT phase-space realization, and GADGET-2 is initialized from the
same CSV fixture (SHA-256 recorded in the GT provenance manifest).  The GADGET
format-1 conversion changes units and stores position and velocity blocks in
single precision; its round-trip audit reports position RMS
$3.43\times10^{-6}$ force-mesh cells and relative momentum RMS
$2.53\times10^{-8}$.  These are small conversion residuals, not exact
bitwise equality of the three solver input files.

All three integrations use the same periodic flat-CDM background, box
$L=168.75\,h^{-1}\,\mathrm{Mpc}$, $128^3$ particles, and the same saved initial
phases at $z_i=99$.  The GADGET-2.0.7 run uses periodic TreePM
\cite{springel2005gadget2} with
$\mathrm{PMGRID}=128$, a short-range tree force, and a halo softening of
$26.367\,h^{-1}\,\mathrm{kpc}$.  Its declared controls include
$\texttt{ErrTolIntAccuracy}=0.025$, $\texttt{MaxRMSDisplacementFac}=0.2$,
$\texttt{MaxSizeTimestep}=0.025$, $\texttt{ErrTolTheta}=0.5$, and
$\texttt{ErrTolForceAcc}=0.005$.  It uses adaptive synchronized timesteps.
IPPL and the native FastPM run instead use plain particle--mesh forces on a
$128^3$ force mesh and the fixed, uniform-in-$\ln a$ schedule described for
ZB.  Thus the underlying collisionless cosmological model and initial phases
are matched, but force resolution, short-range treatment, time stepping, and
input quantization are not identical.  The cross-code difference cannot be
attributed to the TreePM correction alone.

\paragraph{What a Fourier shell means.}
Here ``shell'' means a radial bin in Fourier space, not a layer of particles
in the simulation box.  For integer lattice mode $\boldsymbol n$, the mode
wavenumber is $\boldsymbol k=(2\pi/L)\boldsymbol n$.  The bin index is
$b=\lfloor|\boldsymbol n|+1/2\rfloor$, with $b=1,\ldots,48$; all modes in one
bin have similar $|\boldsymbol k|$.  The plotted center $k_b$ is the
full-lattice-multiplicity-weighted mean $|\boldsymbol k|$ in that bin.
The shell power is the correspondingly weighted mean of the CIC-window-corrected
mode powers, minus Poisson shot noise $1/\bar n=L^3/N_P^3$.  Consequently,
the horizontal axis is $k$ in $h\,\mathrm{Mpc}^{-1}$ and $P(k)$ has units
$(\mathrm{Mpc}/h)^3$.  The first bin is the largest-scale measurement; larger
bin numbers probe progressively smaller scales.

The completed GADGET evolution used eight single-threaded MPI ranks, resumed
from a verified eight-rank checkpoint, and exited normally at $a=1$.  The final
snapshot contains $2{,}097{,}152$ dark-matter particles and records
$z=4.4\times10^{-16}$, numerically zero.  Its measured solver wall time was
2,891.7 seconds.  The same external periodic-CIC estimator was applied to all
three final particle sets on a $128^3$ analysis mesh, including CIC power-window
deconvolution and Poisson shot-noise subtraction.  All 48 bins have positive
subtracted power for each code; they extend through $|\boldsymbol n|=48$, or
75\% of the particle Nyquist frequency.  These upper bins are retained as
characterization, not as a convergence-qualified science band.

\begin{table}[htbp]
\centering
\small
\caption{Selected binwise $z=0$ power differences for the shared $z_i=99$
realization.  Offsets are $100(P_A/P_{\mathrm{FastPM}}-1)$; the final column
uses all 48 positive shot-noise-subtracted bins.  No pass/fail threshold was
set for GT.}
\label{tab:app-gadget2-power}
\setlength{\tabcolsep}{3pt}
\begin{tabular}{@{}>{\raggedright\arraybackslash}p{.38\textwidth}
                >{\centering\arraybackslash}p{.18\textwidth}
                >{\centering\arraybackslash}p{.18\textwidth}
                >{\centering\arraybackslash}p{.20\textwidth}@{}}
\toprule
Pair & Bin 1 ($k=0.0475$) & Bin 48 ($k=1.7873$) & Max. absolute offset \\
 & [$h\,\mathrm{Mpc}^{-1}$] & [$h\,\mathrm{Mpc}^{-1}$] & over 48 bins [\%] \\
\midrule
IPPL PM / FastPM plain PM & $+0.00029\%$ & $+0.00558\%$ & $0.00670\%$ \\
GADGET-2 TreePM / FastPM plain PM & $+0.802\%$ & $+91.241\%$ & $91.241\%$ \\
\bottomrule
\end{tabular}
\end{table}

The IPPL and FastPM spectra nearly overlap over this measured realization.
GADGET-2 is close in the first few bins, then rises above the plain-PM spectrum
toward high $k$; the lower panel reports these offsets relative to FastPM.
This trend is a useful cross-code diagnostic, not evidence that one curve is
the continuum answer: the force models and time integrators differ, there is
only one realization and one resolution, and CIC correction does not remove
finite-mesh aliasing or nonlinear sampling effects.  In particular, the
91\% top-bin difference must not be read as a qualified 91\% force error.

\validationplot{gadget2-zeldovich-comparison.png}{Final $z=0$ matter power for
the shared $z_i=99$ Gaussian 1LPT realization.  Upper panel: shot-noise-
subtracted spectra from IPPL plain PM, native FastPM plain PM, and GADGET-2
TreePM using a common periodic-CIC estimator.  Lower panel: binwise power
offsets relative to FastPM.  A Fourier shell is a radial bin in $|\boldsymbol
n|$, not a configuration-space particle shell.  The GADGET TreePM short-range
force and adaptive synchronized steps differ from the two plain-PM runs.  The
high-$k$ end is characterization, not a convergence claim.  Source: GT.}
{fig:app-gadget2-zeldovich}

\subsection{Evidence integrity, reproducibility, and remaining claims}
\label{app:provenance}

Report identifiers below are relative to
\nolinkurl{build_openmp/demos/cosmology/} in the development worktree, not
archival DOIs.  The \nolinkurl{figures/validation/asset-manifest.json} records
copied plots, original manifests, source-report and asset hashes.  These are
preserved results, not newly executed simulations.

\begingroup
\small
\setlength{\LTcapwidth}{\textwidth}
\begin{longtable}{@{}p{.06\textwidth}>{\raggedright\arraybackslash}p{.89\textwidth}@{}}
\caption{Report identifiers for the completed campaigns.  S uses its released
spatial snapshot; G uses the separately archived Merlin report.}
\label{tab:app-reports}\\
\toprule Key & Authoritative report or exact report snapshot \\
\midrule
\endfirsthead
\toprule Key & Authoritative report or exact report snapshot \\
\midrule
\endhead
L & \nolinkurl{cosmology-validation-na4rs6e0/results.json} \\
ZP & \nolinkurl{zarija-physics-hd3av78e/results.json} \\
ZI & \nolinkurl{zarija-validation-f2c9zb42/results.json} \\
F & \nolinkurl{frozen-force-5_li99qh/results.json} \\
P & \nolinkurl{pancake-validation-rs1o9glb/results.json}; local-gradient audit:
  \nolinkurl{pancake-diagnostics-rs1o9glb/diagnostics.json} \\
E & \nolinkurl{matched-evolution-5v7u3y98/results.json} \\
ER & \nolinkurl{evolution-refinement-k2qeaylz/results.json} \\
S & \nolinkurl{resolution-study-cffjva__/figures-spatial-release/input-report.json.gz} \\
G & \nolinkurl{merlin-cpu-login-20261004/figures-gaussian-release/input-report.json.gz} \\
ZB & \nolinkurl{merlin6:/data/user/adelmann/cosmology-zeldovich-broadband-20261004/analysis/results.json};
  campaign: \nolinkurl{campaign.json}; plot manifest:
  \nolinkurl{plots/manifest.json} \\
GT & \nolinkurl{merlin6:/data/user/adelmann/gadget2-zeldovich-20261004/analysis/gadget2-z99-z0-results.json};
  campaign: \nolinkurl{merlin6:/data/user/adelmann/gadget2-zeldovich-20261004/campaign.json};
  plot manifest: \nolinkurl{plots/three-code-z99-z0-v3/manifest.json} \\
\bottomrule
\end{longtable}
\endgroup

The local spatial stage is complete even though its parent macOS journal later
stopped safely on a disk guard at 44/52 planned runs.  Its ten partial Gaussian
runs are not presented as a completed Gaussian experiment.  G reran all 18
Gaussian cases on Linux with fresh compiler, executable and absolute-path
provenance.  A separate Mac/Linux IC comparison found identical canonical
momenta but position RMS differences about $1.9\times10^{-15}\,h^{-1}\mathrm{Mpc}$;
the CSV bytes are not identical.  This is an IC portability observation, not
a nonlinear cross-platform trajectory qualification.

The Merlin execution checkout remains fixed at \texttt{296fd04c0}.
Documentation/audit handoff is on \texttt{codex/cosmology-linear}, commit
\texttt{9506f1e52}; it must not be confused with the frozen executable source.
Its final integrity audit verifies the canonical 18-run matrix, all 1,355
check identities and qualification flags, 254 comparison records, eight
time-refinement records, 162 numerical archives and 268 file hashes.
The report SHA-256 is\\*
{\footnotesize\texttt{82cb91aae41986c7655a742e71c59150a175f926db1a5e1d2d2709a9ae8156b2}.}

The audit passes while scientific acceptance remains false.  Its role is
completeness, provenance and internal consistency, not an independent
recalculation of all particle physics.  The controller's exit 1 records the
116 scientific failures; there is no execution or disk error in G.
Full tracked-source and executable manifests are also unchanged after plotting.
Numerical NPZ archives preserve exact stored arrays and their digests, but
do not reconstruct the original CSV byte formatting; the earlier gzip-based
E archive does retain original CSV bytes.  Large data remain on Merlin rather
than being duplicated with this manuscript.

ZB is a separate follow-up campaign, not an extension of G's crossed
32/64-resolution matrix.  Its complete analysis report SHA-256 is
listed with the campaign record SHA-256 below.
\begin{quote}\small
Analysis report: \texttt{b345da2fbbf22e9a57091a06352be15d}\\
\texttt{a72d6296809502a571fdc130453f6122}\\
Campaign record: \texttt{340b497e334d50767bf83020185650e3}\\
\texttt{28ec644db498e3b0b2988d7d63ce0129}
\end{quote}
The plot manifest verifies the PNG, SVG, movie, analysis report and campaign
hashes.  A separate copy of the plot manifest is preserved at
\par\smallskip\noindent
\nolinkurl{zeldovich-fastpm-plot-manifest.json}.\par\smallskip
Particle snapshots and the machine-readable spectral arrays remain on Merlin
and are not bundled with this manuscript.

GT is the subsequent single-realization GADGET-2 comparison described in
Section~\ref{app:gadget2-zeldovich}.  The plot-manifest sidecar listed in
Table~\ref{tab:app-reports} records and rechecks the exact report, campaign,
plot-source, PNG, SVG and machine-readable plot-value hashes.  The paper's
copied PNG has the same digest as the rendered figure; the snapshot and full
analysis remain on Merlin.
\begin{quote}\small
Analysis report SHA-256:\\
\texttt{8ad14e25a766665ce3f8c347c250637a}\\
\texttt{2692fd54e06f1a0ea8ac8d3f9b5605a9}\\
Campaign record SHA-256:\\
\texttt{9d7c0fa83907641272c2a7fa677857e4}\\
\texttt{75e897cf858349ea8c7a50fbc35cf80d}\\
Snapshot SHA-256:\\
\texttt{4f10e99fc39146fb5dd17aefc426d07f}\\
\texttt{229df515a4a2dd458eef8345a7df74b5}\\
Paper PNG SHA-256:\\
\texttt{5890ad8f315e9cf2fe0979aa7ac02755}\\
\texttt{a49562a4be338b1a1d74e8747b4948f6}
\end{quote}

Together, the CPU campaigns test background/IC conventions, discrete forces,
bounded trajectories and nonlinear PM comparisons, and MPI reproducibility.
Limits include unresolved local pancake gradients, retained strict mass-sum
failures, and insufficient Gaussian resolution robustness.  Halo statistics,
GPU and multi-node performance remain unqualified.  Larger crossed resolutions,
independent seeds, 2LPT/transient studies, an expanding-universe energy diagnostic,
and backend/scaling tests are subsequent work.  A systematic force- and
step-matched study is still required to isolate the TreePM/plain-PM and
integrator contributions seen in GT.  ZB's very close matched PM spectrum
does not remove any S/G resolution failure or supply a converged science band;
nor does the single-resolution GT comparison qualify a continuum solution.
No failed threshold is retroactively relaxed, and no precision-cosmology or
exascale claim follows from the present CPU campaigns.


\clearpage
\bibliographystyle{unsrt}
\bibliography{references}

\end{document}
